% GENERATED FILE. Edit canonical lessons, then run npm run generate:books.
% canonical-source-hash: fnv1a64:63949554007146f0
% canonical-lessons: IT-C375-annullare, IT-C375-rinviare, IT-C375-provvisorio, IT-C375-definitivo, IT-C375-stabilire

\chapter{Cambio di programma}
\label{ch:it-cambio-di-programma}
\hlchaptermodality{375}

\begin{chapteropening}
\textbf{By the end of this chapter:} \emph{I can change a plan: cancel it, put it off, say it is provisional or final and fix a date.}

Complete the last lesson of chapter 375: I can change a plan: cancel it, put it off, say it is provisional or final and fix a date.
\end{chapteropening}

\section[annullare]{annullare — to cancel, to call off}

\label{lesson:IT-C375-annullare}



\begin{quote}
Before the new one: say the Italian for the alternative, then the Italian for to adapt.
\end{quote}

\subsection*{You'll want to know: annullare}
\textbf{annullare} — \textquotedblleft{}to cancel, to call off\textquotedblright{}.

\begin{quote}
Se piove, annulliamo la gita.
\end{quote}

\begin{quote}
\emph{If it rains, we will call off the trip.}
\end{quote}

\emph{Il volo è stato annullato} — the flight has been cancelled.

\begin{cousinweb}
\emph{annullare} comes from Late Latin \emph{annullāre} \textquotedblleft{}to bring to nothing\textquotedblright{} (EtymDB), built from \emph{ad} and \emph{nūllus} \textquotedblleft{}none\textquotedblright{}.
\end{cousinweb}

\begin{grammarlens}[title={What you've built}]
One more word for telling, explaining and describing at B1.
\end{grammarlens}

\subsection*{Your turn}
\begin{itemize}
\raggedright
  \item \emph{Say it:} \emph{annullare}
  \item \emph{Say it:} \emph{annullare}, once more
  \item \emph{Say it:} a sentence of your own with \emph{annullare}
  \item \emph{Recall:} say the Italian for the alternative, then the Italian for to adapt, then say \emph{annullare} again
\end{itemize}

\begin{tcolorbox}[breakable,colback=teal!4,colframe=teal!35!black,title={Before you move on}]
What does annullare mean? (\textquotedblleft{}To cancel, to call off\textquotedblright{}.) Say it once more.
\end{tcolorbox}
\section[rinviare]{rinviare — to postpone, to put off}

\label{lesson:IT-C375-rinviare}



\begin{quote}
Before the new one: say the Italian for to adapt, then the Italian for to cancel.
\end{quote}

\subsection*{You'll want to know: rinviare}
\textbf{rinviare} — \textquotedblleft{}to postpone, to put off\textquotedblright{}.

\begin{quote}
La riunione è stata rinviata a giovedì.
\end{quote}

\begin{quote}
\emph{The meeting has been put off until Thursday.}
\end{quote}

It also means to send back: \emph{rinviare un pacco}.

\begin{cousinweb}
\emph{rinviare} is built from \emph{ri-} \textquotedblleft{}again, back\textquotedblright{} and \emph{inviare} \textquotedblleft{}to send\textquotedblright{}, usually traced to Late Latin \emph{inviāre} \textquotedblleft{}to put on the way\textquotedblright{}, built on \emph{via} \textquotedblleft{}road\textquotedblright{} (Lewis \& Short).
\end{cousinweb}

\begin{grammarlens}[title={What you've built}]
One more word for telling, explaining and describing at B1.
\end{grammarlens}

\subsection*{Your turn}
\begin{itemize}
\raggedright
  \item \emph{Say it:} \emph{rinviare}
  \item \emph{Say it:} \emph{rinviare}, once more
  \item \emph{Say it:} a sentence of your own with \emph{rinviare}
  \item \emph{Recall:} say the Italian for to adapt, then the Italian for to cancel, then say \emph{rinviare} again
\end{itemize}

\begin{tcolorbox}[breakable,colback=teal!4,colframe=teal!35!black,title={Before you move on}]
What does rinviare mean? (\textquotedblleft{}To postpone, to put off\textquotedblright{}.) Say it once more.
\end{tcolorbox}
\section[provvisorio]{provvisorio — provisional, temporary}

\label{lesson:IT-C375-provvisorio}



\begin{quote}
Before the new one: say the Italian for to cancel, then the Italian for to postpone.
\end{quote}

\subsection*{You'll want to know: provvisorio}
\textbf{provvisorio} — \textquotedblleft{}provisional, temporary\textquotedblright{}.

\begin{quote}
Questo è solo un orario provvisorio.
\end{quote}

\begin{quote}
\emph{This is only a provisional timetable.}
\end{quote}

\emph{In via provvisoria} — for the time being.

\begin{cousinweb}
\emph{provvisorio} probably came from French \emph{provisoire}, from Medieval Latin \emph{prōvīsōrius}, built on \emph{prōvidēre} \textquotedblleft{}to foresee, to provide for\textquotedblright{} (Lewis \& Short), on \emph{vidēre} \textquotedblleft{}to see\textquotedblright{}.
\end{cousinweb}

\begin{grammarlens}[title={What you've built}]
One more word for telling, explaining and describing at B1.
\end{grammarlens}

\subsection*{Your turn}
\begin{itemize}
\raggedright
  \item \emph{Say it:} \emph{provvisorio}
  \item \emph{Say it:} \emph{provvisorio}, once more
  \item \emph{Say it:} a sentence of your own with \emph{provvisorio}
  \item \emph{Recall:} say the Italian for to cancel, then the Italian for to postpone, then say \emph{provvisorio} again
\end{itemize}

\begin{tcolorbox}[breakable,colback=teal!4,colframe=teal!35!black,title={Before you move on}]
What does provvisorio mean? (\textquotedblleft{}Provisional, temporary\textquotedblright{}.) Say it once more.
\end{tcolorbox}
\section[definitivo]{definitivo — final, definitive}

\label{lesson:IT-C375-definitivo}



\begin{quote}
Before the new one: say the Italian for to postpone, then the Italian for provisional.
\end{quote}

\subsection*{You'll want to know: definitivo}
\textbf{definitivo} — \textquotedblleft{}final, definitive\textquotedblright{}.

\begin{quote}
Questa è la data definitiva: non cambia più.
\end{quote}

\begin{quote}
\emph{This is the final date: it will not change again.}
\end{quote}

\emph{In definitiva} — all in all.

\begin{cousinweb}
\emph{definitivo} is a learned borrowing from Latin \emph{dēfīnītīvus} \textquotedblleft{}that defines\textquotedblright{}, from \emph{dēfīnīre} \textquotedblleft{}to set bounds\textquotedblright{} (Lewis \& Short), built on \emph{fīnis} \textquotedblleft{}border, end\textquotedblright{}.
\end{cousinweb}

\begin{grammarlens}[title={What you've built}]
One more word for telling, explaining and describing at B1.
\end{grammarlens}

\subsection*{Your turn}
\begin{itemize}
\raggedright
  \item \emph{Say it:} \emph{definitivo}
  \item \emph{Say it:} \emph{definitivo}, once more
  \item \emph{Say it:} a sentence of your own with \emph{definitivo}
  \item \emph{Recall:} say the Italian for to postpone, then the Italian for provisional, then say \emph{definitivo} again
\end{itemize}

\begin{tcolorbox}[breakable,colback=teal!4,colframe=teal!35!black,title={Before you move on}]
What does definitivo mean? (\textquotedblleft{}Final, definitive\textquotedblright{}.) Say it once more.
\end{tcolorbox}
\section[stabilire]{stabilire — to fix, to set, to establish}

\label{lesson:IT-C375-stabilire}



\begin{quote}
Before the new one: say the Italian for provisional, then the Italian for final.
\end{quote}

\subsection*{You'll want to know: stabilire}
\textbf{stabilire} — \textquotedblleft{}to fix, to set, to establish\textquotedblright{}.

\begin{quote}
Dobbiamo ancora stabilire il giorno della partenza.
\end{quote}

\begin{quote}
\emph{We still have to fix the day of departure.}
\end{quote}

Conjugated with \emph{-isc-}: \emph{stabilisco}. \emph{Stabilirsi} — to settle somewhere.

\begin{cousinweb}
\emph{stabilire} is a learned borrowing from Latin \emph{stabilīre} \textquotedblleft{}to make firm\textquotedblright{} (Lewis \& Short), from \emph{stabilis} \textquotedblleft{}steady\textquotedblright{}, built on \emph{stāre} \textquotedblleft{}to stand\textquotedblright{}.
\end{cousinweb}

\begin{grammarlens}[title={What you've built}]
One more word for telling, explaining and describing at B1.
\end{grammarlens}

\subsection*{Your turn}
\begin{itemize}
\raggedright
  \item \emph{Say it:} \emph{stabilire}
  \item \emph{Say it:} \emph{stabilire}, once more
  \item \emph{Say it:} a sentence of your own with \emph{stabilire}
  \item \emph{Recall:} say the Italian for provisional, then the Italian for final, then say \emph{stabilire} again
\end{itemize}

\begin{tcolorbox}[breakable,colback=teal!4,colframe=teal!35!black,title={Before you move on}]
What does stabilire mean? (\textquotedblleft{}To fix, to set, to establish\textquotedblright{}.) Say it once more.
\end{tcolorbox}
